\ifdefined\gkver\else\def\gkver{student}\fi
\documentclass[\gkver]{gaokaozhenti}
\begin{document}

\chapter{1979年全国}

\begin{enumerate}
\item
%% number: 1
%% paperName: 1979年全国高考第 1 题 3 分
%% typeId: 1
%% type: 单选题
%% chapter: 电路
%% point: 电流强度
%% method: 无
%% score: 3
%% degree: 950
%% duplicateId: 0
%% body:
通过一个电阻的电流是 $5\UA$。经过 $4\Umin$ 时间通过这电阻的一个截面的电量是\xzanswer{C}
\fourchoices[answer=C]
{$20\UC$}
{$50\UC$}
{$1200\UC$}
{$2000\UC$}

%% answer: C
%% memo:
\memoanswer{
由电流的定义式知 $q=It=5\times4\times60=1200\UC$，故 C 正确。
}

\item
%% number: 2
%% paperName: 1979年全国高考第 2 题 3 分
%% typeId: 1
%% type: 单选题
%% chapter: 电路
%% point: 串并联电路
%% method: 无
%% score: 3
%% degree: 850
%% duplicateId: 0
%% body:
把电阻是 $1\UO$ 的一根金属丝截成等长的 10 段，把这 10 段金属丝并联起来。这样并联的一组金属丝的总电阻是\xzanswer{A}
\fourchoices[answer=A]
{$0.01\UO$}
{$0.10\UO$}
{$10\UO$}
{$100\UO$}

%% answer: A
%% memo:
\memoanswer{
金属丝截成 10 段后，每段的电阻为 $r=0.1\UO$，并联后的总电阻为 $R=\frac{r}{10}=0.01\UO$，故 A 正确。
}

\item
%% number: 3
%% paperName: 1979年全国高考第 3 题 3 分
%% typeId: 1
%% type: 单选题
%% chapter: 机械振动
%% point: 单摆
%% method: 无
%% score: 3
%% degree: 750
%% duplicateId: 0
%% body:
单摆的周期在发生下述哪种情况时将增大\xzanswer{D}
\fourchoices[answer=D]
{摆锤的质量增大}
{摆长减小}
{单摆从赤道移到北极}
{单摆从海平面移到高山上}

%% answer: D
%% memo:
\memoanswer{
由单摆周期公式 $T=2\pi\sqrt{\frac{l}{g}}$ 可以看出周期 $T$ 只与 $l$ 及 $g$ 有关，与摆锤质量无关，故 A 错误；摆长减小，周期减小，故 B 错误；从赤道到北极，$g$ 增大，所以周期减小，故 C 错误；从海平面到高山，$g$ 减小，所以周期增大，故 D 正确。
}

\item
%% number: 4
%% paperName: 1979年全国高考第 4 题 3 分
%% typeId: 1
%% type: 单选题
%% chapter: 力物体的平衡
%% point: 受力分析
%% method: 无
%% score: 3
%% degree: 850
%% duplicateId: 0
%% body:
放在光滑斜面上加速下滑的物体受到的力是\xzanswer{A}
\fourchoices[answer=A]
{重力和斜面支持力}
{重力、下滑力和斜面支持力}
{重力、斜面支持力和加速力}
{重力、斜面支持力、下滑力和正压力}

%% answer: A
%% memo:
\memoanswer{
光滑的斜面上加速下滑的物体受到重力和支持力的作用，故 A 正确；重力沿斜面的分力使物体沿斜面加速下滑，而不存在一个独立于重力之外的所谓“下滑力”或者“加速力”，故 BC 错误；压力作用在斜面上，不是物体所受到的力，故 D 错误。
}

\item
%% number: 5
%% paperName: 1979年全国高考第 5 题 3 分
%% typeId: 1
%% type: 单选题
%% chapter: 光学
%% point: 光的折射
%% method: 无
%% score: 3
%% degree: 750
%% duplicateId: 0
%% body:
光线透过空气中的平行平面厚玻璃板，问下图所示四种情形中哪一种是正确的\xzanswer{B}
\fourpicture[w=3.1cm]{figs/05a.png}{figs/05b.png}{figs/05c.png}{figs/05d.png}
\fourchoices[answer=B]
{$i=i'$}
{$i=i'$，$i>r$}
{$i=i'$，$i<r$}
{$i>i'$，$i>r$}

%% answer: B
%% memo:
\memoanswer{
由空气倾斜进入玻璃板，折射角小于入射角，即 $i>r$，故 AC 错误；当光再由玻璃板斜射进入空气时，由几何关系可得光在下表面的入射角等于上表面的折射角，故从下表面射出的折射角等于上表面的入射角，即 $i=i'$，出射光线与最初的入射光线平行，故 D 错误 B 正确。
}

\item
%% number: 6
%% paperName: 1979年全国高考第 6 题 3 分
%% typeId: 1
%% type: 单选题
%% chapter: 原子和原子核
%% point: 人工核反应
%% method: 无
%% score: 3
%% degree: 850
%% duplicateId: 0
%% body:
中子（\ce{^{1}_{0}n}）轰击硼核（\ce{^{10}_{5}B}）发生核反应如下：\ce{^{10}_{5}B + ^{1}_{0}n -> ^{7}_{3}Li} + （  ），括号中是\xzanswer{A}
\fourchoices[answer=A]
{$\alpha$ 粒子}
{质子}
{中子}
{$\beta$ 粒子}

%% answer: A
%% memo:
\memoanswer{
由质量数和核电荷数守恒可知放出的粒子质量数为 4，电荷数为 2，故该粒子为 $\ce{^{4}_{2}He}$，即 $\alpha$ 粒子，故 A 正确。
}

\item
%% number: 7
%% paperName: 1979年全国高考第 7 题 3 分
%% typeId: 1
%% type: 单选题
%% chapter: 气体性质
%% point: 物体的内能
%% method: 无
%% score: 3
%% degree: 750
%% duplicateId: 0
%% body:
把 $100\Ug$ $0\UdC$ 的冰投入 $100\Ug$ $50\UdC$ 的水中，混合时与外界无热量交换，达到热平衡后的温度是\xzanswer{C}
\fourchoices[answer=C]
{$5\UdC$}
{$-15\UdC$}
{$0\UdC$}
{$4\UdC$}

%% answer: C
%% memo:
\memoanswer{
水的相变潜热为 $335\UkJkg = 335\UJg$，$100\Ug$ $0\UdC$ 的冰熔化为 $0\UdC$ 水吸热 $335\UJg \times 100\Ug = 33500\UJ$，而 $100\Ug$ $50\UdC$ 水降到 $0\UdC$ 放热 $4.2\UJgdC \times 100\Ug \times 50\UdC = 21000\UJ$，显然吸热量比放热量大，就说明冰还没有完全熔化水就降到 $0\UdC$，温度相同就没有热传递，最后得到的是冰水混合物，温度为 $0\UdC$，故 C 正确。
}

\item
%% number: 8
%% paperName: 1979年全国高考第 8 题 3 分
%% typeId: 1
%% type: 单选题
%% chapter: 电场
%% point: 电容
%% method: 无
%% score: 3
%% degree: 750
%% duplicateId: 0
%% body:
一个平行板电容器充电后，把电源断开，再用绝缘的工具把电容器的两金属板拉开一些，这使\xzanswer{D}
\fourchoices[answer=D]
{电容器中的电量增加}
{电容器的电容增加}
{电容器的电压不变}
{电容器的电压增加}

%% answer: D
%% memo:
\memoanswer{
由电容的决定式 $C=\frac{\varepsilon S}{4\pi kd}$ 知，电容与板间距离成反比，当把两板间距拉大，电容减小，故 B 错误；电容器充电后断开开关，电容器所带电荷量保持不变，故 A 错误；由 $Q=CU$ 知电压增大，故 C 错误 D 正确。
}

\item
%% number: 9
%% paperName: 1979年全国高考第 9 题 3 分
%% typeId: 1
%% type: 单选题
%% chapter: 运动和力
%% point: 牛顿第二定律
%% method: 无
%% score: 3
%% degree: 650
%% duplicateId: 0
%% body:
一条绳能承受的最大拉力是 $100\UN$（超出此值，绳就被拉断）。用这条绳拉一个质量是 $2\Ukg$ 的物体并使其仍在光滑的水平面上运动，绳和水平面的夹角是 $60^{\circ}$。在绳不被拉断的条件下，物体的最大加速度可以达到\xzanswer{A}
\fourchoices[answer=A]
{$5.77\Umsq$}
{$25\Umsq$}
{$43\Umsq$}
{$100\Umsq$}

%% answer: A
%% memo:
\memoanswer{
\onepicture[w=4.5cm]{figs/09_an.png}

如图所示，若保证物体始终不离开地面，设施加的绳子拉力最大为 $F$，此时水平面对物体的支持力为 0，竖直方向有 $F\sin60^{\circ}=mg$，得 $F=\frac{40\sqrt{3}}{3}\UN<100\UN$，在绳子承受的拉力范围内，此时物体的加速度 $a$ 可达到最大，水平方向有 $F\cos60^{\circ}=ma$，代入数据得 $a\approx5.77\Umsq$，故 A 正确。
}

\item
%% number: 10
%% paperName: 1979年全国高考第 10 题 3 分
%% typeId: 1
%% type: 单选题
%% chapter: 功和能
%% point: 功率
%% method: 无
%% score: 3
%% degree: 750
%% duplicateId: 0
%% body:
质量是 5 吨的汽车在水平路面上以加速度 $a=2\Umsq$ 起动，所受阻力是 $1.0\times10^{3}\UN$。汽车起动后第一秒末的瞬时功率是\xzanswer{D}
\fourchoices[answer=D]
{$2\UkW$}
{$11\UkW$}
{$20\UkW$}
{$22\UkW$}

%% answer: D
%% memo:
\memoanswer{
汽车在启动时水平方向受到牵引力和摩擦力的作用，由牛顿第二定律得 $F-f=ma$，解得 $F=f+ma=1000\UN+5000\times2\UN=11000\UN$，由速度公式知汽车第 $1\Us$ 末的速度为 $v=at=2\times1\Ums=2\Ums$，所以 $P=Fv=11000\times2\UW=22000\UW=22\UkW$，故 D 正确。
}

\item
%% number: 11
%% paperName: 1979年全国高考第 11 题 3 分
%% typeId: 1
%% type: 单选题
%% chapter: 磁场
%% point: 洛伦兹力
%% method: 无
%% score: 3
%% degree: 650
%% duplicateId: 0
%% body:
一个长螺线管通有交流电，把一个带电粒子沿管轴射入管中，粒子将在管中\xzanswer{D}
\fourchoices[answer=D]
{做圆周运动}
{沿轴线来回运动}
{做匀加速直线运动}
{做匀速直线运动}

%% answer: D
%% memo:
\memoanswer{
螺线管中通入交流电后，将在管中产生交变的磁场，但磁场方向总是与轴线平行，带电粒子不受洛伦兹力的作用，因此粒子将做匀速直线运动，故 D 正确。
}

\item
%% number: 12
%% paperName: 1979年全国高考第 12 题 10 分
%% typeId: 6
%% type: 计算题
%% chapter: 光学
%% point: 透镜
%% method: 无
%% score: 10
%% degree: 650
%% duplicateId: 0
%% body:
一架显微镜物镜的焦距 $f_{1}=0.3\Ucm$，目镜的焦距 $f_{2}=2\Ucm$，目镜和物镜相距 $16\Ucm$。如果从这显微镜观察到的像离目镜 $25\Ucm$，被观察的物体离物镜多远？

%% answer: 0.31 cm
\jdanswer{
$0.31\Ucm$
}

%% memo:
\memoanswer{
从显微镜的工作原理可知，物镜成像为实像、目镜成像应为虚像，且两镜均为凸透镜，成像示意图如图所示。

\onepicture[h=5.5cm]{\includesvg{figs/12_an}}

$A_{1}B_{1}$ 经过目镜成虚像于 $A_{2}B_{2}$，由成像公式得

$\frac{1}{u_{2}}+\frac{1}{v_{2}}=\frac{1}{f_{2}}$

代入数据 $v_{2}=-25\Ucm$，$f_{2}=2\Ucm$，得

$\frac{1}{u_{2}}=\frac{1}{f_{2}}-\frac{1}{v_{2}}=\frac{1}{2}+\frac{1}{25}=\frac{27}{50}$

故目镜的物距 $u_{2}=\frac{50}{27}\Ucm$

物体 AB 通过物镜成实像于 $A_{1}B_{1}$，则

$v_{1}=16-\frac{50}{27}\approx14.15\Ucm$

由成像公式得

$\frac{1}{u_{1}}+\frac{1}{v_{1}}=\frac{1}{f_{1}}$

解得 $\frac{1}{u_{1}}=\frac{1}{f_{1}}-\frac{1}{v_{1}}=\frac{1}{0.3}-\frac{1}{14.15}$

故物体 AB 距物镜的距离 $u_{1}\approx0.31\Ucm$。

评分标准：全题 10 分，（1）6 分，（2）3 分，（3）1 分。

（1）能正确利用公式计算目镜物距，给 6 分。只正确写出透镜公式，不知道目镜的像是虚像，须用负数代入的，只给 2 分。如果说明是虚像，但未用负数代入的，给 3 分。透镜公式写错但知道目镜所成的像是虚像，像距用负数计算的，给 2 分。

（2）求出物镜的像距，给 3 分。不会用显微镜结构来计算物镜的像距，不给分。

（3）求出物镜的物距，给 1 分。
}

\item
%% number: 13
%% paperName: 1979年全国高考第 13 题 10 分
%% typeId: 6
%% type: 计算题
%% chapter: 气体性质
%% point: 玻意耳定律
%% method: 无
%% score: 10
%% degree: 450
%% duplicateId: 0
%% body:
一个潜水艇位于水面下 $200\Um$。艇上有一个容积 $V_{1}=2\Umc$ 的贮气钢筒，筒内贮有压缩空气。将筒内一部分空气压入水箱（水箱有排水孔和海水相连）排出海水 $10\Umc$。此时筒内剩余气体的压强是 95 个大气压。设在排水过程中温度不变，求贮气钢筒内原来压缩空气的压强。（计算时取 1 个大气压 $p_{0}=1.0\Ukgcmq$，海水密度取 $1.0\times10^{3}\Ukgmc$）

%% answer: 200 个大气压
\jdanswer{
200 个大气压
}

%% memo:
\memoanswer{
令 $p_{0}=1$ 个大气压。

解法一：筒内原有压缩空气物质的量为 $n$，压强为 $p_{1}$，体积为 $V_{1}=2\Umc$，由理想气体状态方程得

$p_{1}V_{1}=nRT$，解得 $n=\frac{p_{1}\times2}{RT}$，

排水后筒内压缩空气物质的量为 $n_{1}$，压强为 $p_{1}'=95p_{0}$，体积仍为 $V_{1}=2\Umc$，由理想气体状态方程得

$p_{1}'V_{1}=n_{1}RT$，解得 $n_{1}=\frac{p_{1}'V_{1}}{RT}=\frac{95p_{0}\times2}{RT}=\frac{190p_{0}}{RT}$，

水箱中空气物质的量为 $n_{2}$，体积 $V_{2}=10\Umc$，压强约等于 200 米深处水的压强，水产生的压强为 $p_{\text{水}}$，有

$p_{\text{水}}=\rho gh=2.0\times10^{6}\UPa$，

又 $p_{0}=1.0\times10^{5}\UPa$，

故 $p_{\text{水}}=20p_{0}$，

由压强平衡得

$p_{2}=p_{\text{水}}+p_{0}=21p_{0}$，

由理想气体状态方程得

$p_{2}V_{2}=n_{2}RT$，解得 $n_{2}=\frac{p_{2}V_{2}}{RT}=\frac{21p_{0}\times10}{RT}=\frac{210p_{0}}{RT}$，

又 $n=n_{1}+n_{2}$，

即 $\frac{p_{1}\times2}{RT}=\frac{190p_{0}}{RT}+\frac{210p_{0}}{RT}$，解得 $p_{1}=200p_{0}$，

所以筒内原有气体的压强为 200 大气压。

解法二：贮气筒原来压强为 $p_{1}$，体积为 $V_{1}=2\Umc$，假如筒内气体等温膨胀至压强 $p_{1}'=95p_{0}$，则此状态下气体体积为

$V_{1}'=\frac{p_{1}V_{1}}{p_{1}'}=\frac{2}{95}p_{1}(\Umc)$，

因为贮气筒已占去了压强为 $95p_{0}$ 的气体 $2\Umc$，故潜水艇水箱内的气体，可以看做是压强为 $p_{1}'=95p_{0}$，体积为 $(V_{1}'-2)\Umc$ 的气体被等温压入至水箱内，压入水箱后气体的体积 $V_{2}=10\Umc$，压强 $p_{2}=\rho gh+p_{0}=21p_{0}$，由玻意耳定律得

$p_{1}'(V_{1}'-2)=p_{2}V_{2}$，

即 $95\left(\frac{2}{95}p_{1}-2\right)=21p_{0}\times10$，

解得 $p_{1}=200p_{0}$，

所以贮气筒内原有气体压强为 200 大气压。

解法三：若贮气筒原来压强为 $p_{1}$，体积为 $V_{1}=2\Umc$，因为气体处于均匀状态，可以设想为其中有 $\frac{V_{1}}{n}$ 的气体被压入潜水艇水箱内，压强变为 $p_{2}=p_{0}+\rho gh=21p_{0}$，体积为 $V_{2}=10\Umc$，筒内余下 $\left(1-\frac{1}{n}\right)V_{1}$ 的气体，后来充满了 $V_{1}$，且压强变为 $p_{1}'=95p_{0}$，由玻意耳定律得

$p_{1}\frac{V_{1}}{n}=p_{2}V_{2}$，

$p_{1}\left(1-\frac{1}{n}\right)V_{1}=p_{1}'V_{1}$，

把上两式相加得

$p_{1}V_{1}=p_{1}'V_{1}+p_{2}V_{2}$，

即 $p_{1}\times2=95p_{0}\times2+21p_{0}\times10$，

解得 $p_{1}=200p_{0}$，

所以筒内原来气体的压强为 200 个大气压。

评分标准：全题 10 分。（1）4 分，（2）6 分。

（1）正确算出水箱中空气压强的，给 4 分。漏算海面大气压的，扣 1 分。只列出 $\langle1\rangle$，没有代入数据计算的，只给 2 分。

$\langle1\rangle$ 式中没有写 $g$，但计算正确的，不扣分。单位不配套的，扣 1 分。

（2）正确算出贮气筒中原有气体压强的，给 6 分。

只列出 $\langle2\rangle$ 式、$\langle3\rangle$ 式的，各给 1 分。

只导出或只直接写出 $\langle4\rangle$ 式的，给 4 分。
}

\item
%% number: 14
%% paperName: 1979年全国高考第 14 题 13 分
%% typeId: 6
%% type: 计算题
%% chapter: 曲线运动
%% point: 竖直平面圆周运动
%% method: 无
%% score: 13
%% degree: 450
%% duplicateId: 0
%% body:
有一质量是 $m$ 的小球 Ⅰ 用长度是 $l$ 的绳子悬挂在 O 点。把球 Ⅰ 拉到 A 点，OA 是水平线，如图所示，另一质量相等的小球 Ⅱ 静止放在 B 点（B 点在 O 点的竖直下方，$OB=l$），BC 是半径 $R=\frac{l}{2}$ 的一段圆弧轨道（圆心在 OB 的中点）。当在 A 点的小球 Ⅰ 从静止下落到 B 点时，跟小球 Ⅱ 作弹性碰撞使小球 Ⅱ 沿轨道 BC 滑出（不考虑摩擦）。求小球 Ⅱ 经过 D 点时对轨道的压力。（圆弧轨道 BD 所对的圆心角 $\theta=60^{\circ}$，$m=1\Ukg$，$g$ 用 $10\Umsq$ 计算）
\onepicture[align=right, w=7cm]{figs/14.png}

%% answer: 35 N
\jdanswer{
$35\UN$
}

%% memo:
\memoanswer{
球 Ⅰ 到达 B 点时的速度 $v_{1}$，由机械能守恒得

$mgl=\frac{1}{2}mv_{1}^{2}$，解得 $v_{1}=\sqrt{2gl}$

Ⅰ、Ⅱ 两球在 B 点作弹性碰撞，球 Ⅰ 碰撞后的速度 $v_{1}'$，球 Ⅱ 碰撞后的速度 $v_{\mathrm{II}}$，由动量守恒和机械能守恒得

$mv_{\mathrm{I}}=mv_{\mathrm{I}}'+mv_{\mathrm{II}}$，$\frac{1}{2}mv_{\mathrm{I}}^{2}=\frac{1}{2}mv_{\mathrm{I}}'^{2}+\frac{1}{2}mv_{\mathrm{II}}^{2}$，

联立解得 $v_{\mathrm{II}}=\sqrt{2gl}$，$v_{\mathrm{I}}'=0$，

球 Ⅱ 上滑到 D 点时的速度 $v_{\mathrm{II}}'$，由机械能守恒得

$\frac{1}{2}mv_{\mathrm{II}}^{2}=\frac{1}{2}mv_{\mathrm{II}}'^{2}+mgR(1-\cos\theta)$

解得 $v_{\mathrm{II}}'^{2}=\frac{3}{2}gl$ ①，

设 $N$ 是在 D 点轨道对球 Ⅱ 的支持力，由牛顿第二定律得

$N-mg\cos\theta=\frac{mv_{\mathrm{II}}'^{2}}{R}$ ②

把①和 $R=\frac{l}{2}$、$\cos\theta=\frac{1}{2}$、$m=1\Ukg$ 代入②得

$N=mg\cos\theta+m\frac{v_{\mathrm{II}}'^{2}}{R}=35\UN$。

球 Ⅱ 对轨道的压力 $N'$ 同轨道对球 Ⅱ 的压力 $N$ 大小相等方向相反，故 $N'=N=35\UN$。

评分标准：全题 13 分。（1）2 分，（2）3 分，（3）3 分，（4）5 分。

（1）知道利用能量守恒并写出公式 $\langle1\rangle$，给 2 分。

（2）知道 $m_{\mathrm{I}}=m_{\mathrm{II}}$ 的两球作弹性碰撞时，$v_{\mathrm{I}}=v_{\mathrm{II}}$，并求出 $v_{\mathrm{II}}$ 的表达式 $\langle2\rangle$ 的，给 3 分，根据能量守恒和动量守恒定律只写出公式而未算出 $v_{\mathrm{II}}$ 的表达式的，只给 2 分。

（3）知道利用能量守恒关系式导出 $v_{\mathrm{II}}'$ 的表达式 $\langle3\rangle$，给 3 分。

（4）知道求压力时考虑向心力和重力，得出正确答案，给 5 分。未说明轨道对球的压力和球对轨道压力的关系的，不扣分。

如直接从小球 Ⅰ 的位能写出 $v_{\mathrm{II}}'$，中间过程未作文字说明的，扣去 $\langle2\rangle$ 中的 3 分，如只未说明速度交换的，扣 1 分；如对能量守恒，动量守恒和速度交换作了文字说明，虽未写出中间过程的公式，不扣分。

在分析压力时，知道要考虑向心力，但方向分析错误，或知道向心力，但对重力未分解或分析错误，扣 3 分。
}

\item
%% number: 15
%% paperName: 1979年全国高考第 15 题 13 分
%% typeId: 6
%% type: 计算题
%% chapter: 电路
%% point: 闭合电路欧姆定律
%% method: 无
%% score: 13
%% degree: 450
%% duplicateId: 0
%% body:
有电路如图，$R_{1}=3000\UO$，$V_{\mathrm{A}}$ 是内阻为 $6000\UO$ 的电压表，$V_{\mathrm{B}}$ 是内阻为 $3000\UO$ 的电压表。已知：

$K_{1}$ 断开，$K_{2}$ 接到 A 时，电压表的读数是 $4\UV$；

$K_{1}$ 接通，$K_{2}$ 接到 A 时，电压表的读数是 $8\UV$；

$K_{1}$ 接通，$K_{2}$ 接到 B 时，电压表的读数是 $7.5\UV$。

求 $R_{2}$ 的值。
\onepicture[align=right, w=7cm]{figs/15.png}

%% answer: 2500 Ω
\jdanswer{
$2500\UO$
}

%% memo:
\memoanswer{
由题意知，$R_{1}$ 与 $V_{\mathrm{A}}$ 并联的总电阻为

$R_{1\mathrm{A}}=\frac{R_{1}R_{\mathrm{A}}}{R_{1}+R_{\mathrm{A}}}=\frac{3000\times6000}{3000+6000}\UO=2000\UO$；

$R_{1}$ 与 $V_{\mathrm{B}}$ 并联的总电阻为

$R_{1\mathrm{B}}=\frac{R_{1}R_{\mathrm{B}}}{R_{1}+R_{\mathrm{B}}}=\frac{3000}{2}\UO=1500\UO$。

设电源电动势是 $E$，内阻是 $r$，由闭合电路欧姆定律，对给出的三种情况列出三个方程：当 $K_{1}$ 断开，$K_{2}$ 接到 A 时，电压表示数 $U_{1}=4\UV$，有

$E=U_{1}+\frac{U_{1}}{R_{1\mathrm{A}}}(R_{2}+R_{3}+r)$ ①，

当 $K_{1}$ 接通，$K_{2}$ 接到 A 时，电压表示数 $U_{2}=8\UV$，有

$E=U_{2}+\frac{U_{2}}{R_{1\mathrm{A}}}(R_{3}+r)$ ②，

当 $K_{1}$ 接通，$K_{2}$ 接到 B 时，电压表示数 $U_{3}=7.5\UV$，有

$E=U_{3}+\frac{U_{3}}{R_{1\mathrm{B}}}(R_{3}+r)$ ③，

令 $R_{3}+r=R$，解②③得

$E=10\UV$，$R=500\UO$，

代入①得 $R_{2}=2500\UO$。

评分标准：全题 13 分。（1）1 分，（2）1 分，（3）9 分，（4）2 分。

（1）计算 $R_{1}$ 与 $V_{\mathrm{A}}$ 并联的总电阻，1 分。公式错、运算错不给分，不写单位不扣分。

（2）计算 $R_{1}$ 与 $V_{\mathrm{B}}$ 并联的总电阻，1 分。公式错、运算错不给分，不写单位不扣分。

（3）正确列出方程 $\langle1\rangle$、$\langle2\rangle$、$\langle3\rangle$，给 9 分（每方程 3 分）。

如果不把电压表的内阻计入电路参数，虽然形式列出三个方程，仍扣 3 分。（因为不合物理原理）。

（4）算出 $R_{2}$，2 分。

如果没有提到 $r$ 或设 $r$ 为零解题，同样给分。
}

\item
%% number: 16
%% paperName: 1979年全国高考第 16 题 13 分
%% typeId: 6
%% type: 计算题
%% chapter: 运动和力
%% point: 牛顿第二定律
%% method: 无
%% score: 13
%% degree: 550
%% duplicateId: 0
%% body:
在加速行驶的火车上固定一斜面，斜面角是 $\theta$。有一物体静止在斜面上。如果火车加速度小于某一值 $a_{0}$，物体就会下滑。设物体和斜面间的静摩擦系数是 $\mu$，推导 $a_{0}$ 的表达式。（静摩擦系数 = $\frac{\text{最大静摩擦力}}{\text{正压力}}$）

%% answer: $a_{0} = \frac{\sin \theta  - {\mu}\cos \theta }{{\cos \theta  + {\mu}\sin \theta }} g$
\jdanswer{
$a_{0}=\frac{\sin\theta-\mu\cos\theta}{\cos\theta+\mu\sin\theta}g$
}

%% memo:
\memoanswer{
解法一：如图 1 作一正交坐标系，在地面观察者看来物体以加速度 $a_{0}$ 沿 $x$ 轴方向前进，由受力分析和牛顿第二定律得

在 $y$ 轴方向上：$N\cos\theta+f_{0}\sin\theta=mg$ ①

在 $x$ 轴方向上：$N\sin\theta-f_{0}\cos\theta=ma_{0}$ ②

又 $f_{0}=\mu N$ ③

由②式除以①式得 $\frac{a_{0}}{g}=\frac{N\sin\theta-\mu N\cos\theta}{N\cos\theta+\mu N\sin\theta}$，

解得 $a_{0}=\frac{\sin\theta-\mu\cos\theta}{\cos\theta+\mu\sin\theta}g$

\onepicture[num=1, h=3.5cm]{figs/16_an1.png}

解法二：如图 2 作一正交坐标系，地面观察者看来物体以加速度 $a_{0}$ 前进，把 $a_{0}$ 沿两坐标轴分解得

$a_{x}=a_{0}\cos\theta$，$a_{y}=a_{0}\sin\theta$

由受力分析和牛顿第二定律得

在 $x$ 轴方向上：$mg\sin\theta-f_{0}=ma_{0}\cos\theta$ ①

在 $y$ 轴方向上：$N-mg\cos\theta=ma_{0}\sin\theta$ ②

又 $f_{0}=\mu N$ ③

\onepicture[num=2, h=3.5cm]{figs/16_an2.png}

联立解得 $a_{0}=\frac{\sin\theta-\mu\cos\theta}{\cos\theta+\mu\sin\theta}g$

评分标准：全题 13 分。①5 分，②5 分，③3 分。

（1）正确写出①式，给 5 分。

①式有错误，但能正确写出这三项的大小的，给 2 分。漏掉一项的，不给分。

（2）正确写出②式，给 5 分。

②式有错误，但能正确写出这三项的大小的，给 2 分。漏掉一项的，不给分。

（3）由①、②式解出③式的，给 3 分。答案中还出现正压力的，不给分。

按其他方向分解，例如在竖直和水平方向进行分解的，参照上述标准给分。用惯性力概念分析（即把 $ma$ 方向反转，然后按静力学的方法计算），得到结果的也参照上述标准给分。
}

\item
%% number: 17
%% paperName: 1979年全国高考第 17 题 4 分
%% typeId: 1
%% type: 单选题
%% chapter: 磁场
%% point: 左手定则
%% method: 无
%% score: 4
%% degree: 450
%% duplicateId: 0
%% body:
两条直导线互相垂直（如图），但相隔一个小的距离，其中一条 AB 是固定的，另一条 CD 能自由活动。当直流电流按图中所示方向通入两条导线时，导线 CD 将\xzanswer{E}
\onepicture[w=6cm]{figs/17.png}
\fivechoices[answer=E]
{不动}
{顺时针方向转动，同时靠近导线 AB}
{逆时针方向转动，同时离开导线 AB}
{顺时针方向转动，同时离开导线 AB}
{逆时针方向转动，同时靠近导线 AB}

%% answer: E
%% memo:
\memoanswer{
电流 AB 产生的磁场在右边垂直纸面向里，在左边垂直纸面向外，在 CD 左右两边各取一小电流元，由左手定则知，左边的电流元所受的安培力方向向下，右边的电流元所受安培力方向向上，所以 CD 导线逆时针方向转动。当 CD 导线转过 $90^{\circ}$ 后，两电流为同向电流，相互吸引，所以导线 CD 逆时针方向转动的同时靠近导线 AB，故 E 正确。
}

\item
%% number: 18
%% paperName: 1979年全国高考第 18 题 4 分
%% typeId: 1
%% type: 单选题
%% chapter: 电磁感应
%% point: 法拉第电磁感应定律
%% method: 无
%% score: 4
%% degree: 450
%% duplicateId: 0
%% body:
一环形闭合线圈放在均匀磁场中，线圈的轴线与磁场方向成 $30^{\circ}$ 角，磁感应强度随时间均匀变化。在下述办法中（如需要改绕线圈，用原规格的导线），用哪一种办法可以使线圈中的感生电流增加一倍\xzanswer{C}
\fivechoices[answer=C]
{把线圈的匝数增加一倍}
{把线圈的面积增加一倍}
{把线圈的半径增加一倍}
{改变线圈轴线对磁场的方向}
{把线圈的匝数减少到原来的一半}

%% answer: C
%% memo:
\memoanswer{
设感应电流为 $I$，电阻为 $R$，匝数为 $n$，线圈半径为 $r$，线圈面积为 $S$，由法拉第电磁感应定律得 $E=n\frac{\Delta\phi}{\Delta t}=n\frac{\Delta B\cdot S\cos30^{\circ}}{\Delta t}$，由欧姆定律得 $I=\frac{E}{R}$，线圈导线总长度 $l=n\cdot2\pi r$，$S=\pi r^{2}$，电阻 $R=\rho\frac{l}{S'}$（$S'$ 为导线截面积），则 $I=\frac{\sqrt{3}S'}{4\rho}\cdot\frac{\Delta B}{\Delta t}\cdot r$，其中 $\frac{\Delta B}{\Delta t}$ 是磁感应强度的变化率，为恒量，电阻率 $\rho$ 和导线截面积 $S'$ 也是恒量，所以 $I\propto r$，故 C 正确；线圈的匝数增加 1 倍或减少一半时，感生电流的大小不变，故 AE 错误；线圈的面积增加 1 倍，线圈的半径 $r$ 增大 $\sqrt{2}$ 倍，感生电流的增大 $\sqrt{2}$ 倍，故 B 错误；感应电流的大小与转轴对磁场的方向无关，故 D 错误。

评分标准：全题 8 分；每小题正确答案 4 分，错误答案 $-1$ 分，不答的，不给分。
}

\end{enumerate}

\end{document}
